% Numeração das subseções como letras
\renewcommand{\thesubsection}{\Alph{subsection}}

Nesta seção, são descritos os métodos que orientaram as atividades desta iniciação científica. São apresentados aspectos relacionados à classificação da pesquisa, seguidos do levantamento bibliográfico e dos critérios de seleção do referencial teórico.

\subsection{Classificação da Pesquisa}
\label{ssec:classificacao-pesquisa}

Este estudo é caracterizado como uma pesquisa aplicada, de abordagem quantitativa e natureza experimental. Quanto aos objetivos, classifica-se como exploratória e explicativa, buscando compreender e explicar fenômenos relacionados à aplicação de IA no diagnóstico oftalmológico. Os procedimentos adotados incluem levantamento bibliográfico, desenvolvimento e validação de modelos computacionais, em contextos transversais e longitudinais.

\paragraph{Técnicas metodológicas empregadas}
\begin{itemize}
  \item Extração e pré-processamento de imagens retinográficas;
  \item Desenvolvimento e treinamento de modelos baseados em clustering, SVM e CNN;
  \item Avaliação comparativa de desempenho entre modelos treinados a partir de bases locais e públicas;
  \item Análise exploratória para identificar características relevantes das imagens associadas às condições oftalmológicas;
  \item Validação cruzada e testes de robustez dos modelos para garantir a aplicabilidade prática dos resultados.
\end{itemize}

\subsection{Levantamento Bibliográfico}
\label{ssec:levantamento-bibliografico}

O levantamento bibliográfico foi realizado nas bases PubMed, IEEE Xplore, Web of Science e Scopus, utilizando as seguintes strings de busca refinadas conforme a compatibilidade dos resultados com o escopo da pesquisa. A Tabela~\ref{tab:grade-completa} apresenta as consultas originais, as bases consultadas e o número de artigos obtidos, sendo filtrados apenas os resultados de publicações dos últimos dez anos (2015-2025).

\begin{table}[ht]
  \centering
  \caption{Strings de busca, bases de dados consultadas e número de resultados encontrados}
  \label{tab:grade-completa}
  \begin{tabular}{@{} l c r @{}}
    \toprule
    \textbf{String de Busca} & \textbf{Base de Dados} & \textbf{Resultados} \\ 
    \midrule
    \texttt{"diabetic retinopathy"}                     & Web of Science & 22\,528 \\ 
    \texttt{"diabetic retinopathy"}                     & Scopus         & 32\,473 \\ 
    \texttt{"diabetic retinopathy"}                     & PubMed         &  3\,542 \\ 
    \texttt{"diabetic retinopathy" AND brazil}          & Web of Science &     10 \\ 
    \texttt{"diabetic retinopathy" AND brazil}          & Scopus         &     11 \\ 
    \bottomrule
  \end{tabular}
\end{table}
\begin{flushleft}
  \small Fonte: Autor.
\end{flushleft}

\subsection{Critérios de Seleção}
\label{ssec:criterios-selecao}

Os artigos selecionados atenderam aos seguintes critérios:
\begin{itemize}
  \item Relevância científica (contribuições teóricas ou empíricas significativas);
  \item Impacto (fator de impacto dos periódicos e número de citações);
  \item Publicação recente (últimos dez anos, 2015-2025);
  \item Pertinência metodológica com o tema da pesquisa.
\end{itemize}